% Evaluation, organised by research question (RQ list, one subsection per RQ, each closed by
% a shaded summary box), in the style of Ghosh et al., arXiv:2604.19855.
% Numbers: test/runs/2026-10-30-real13 (143 runs), every run at its own code distance
% (fit_d.py, p' = p). Figures: test/runs/2026-10-30-real13/analysis/figs_v2/.
%
% Preamble:
%   \usepackage{enumitem}
%   \usepackage[most]{tcolorbox}
%   \newtcolorbox{rqsummary}{enhanced, colback=blue!7, colframe=blue!7, boxrule=0pt,
%     sharp corners, left=4pt, right=4pt, top=3pt, bottom=3pt,
%     before skip=6pt, after skip=6pt}

% ---- end of the evaluation setup, after the metrics ---------------------------------------

Our evaluation answers three research questions:
\begin{description}[leftmargin=2.6em, labelwidth=2.2em, font=\normalfont\bfseries]
  \item[RQ1:] With the same kind of T-state source and the same error target, does moving
    the sources from the data plane to an upper layer lower space-time volume?
  \item[RQ2:] Does the reduction come from fewer physical qubits or from fewer rounds?
  \item[RQ3:] How many upper-layer sources does HBM need?
\end{description}
% RQ4 (cost of a cross-layer merge) once the seam simulation is in.

% ---- RQ1 ------------------------------------------------------------------------------------

\subsection{(RQ1) Space-time volume}

\textbf{Like-for-like comparison:} Fig.~\ref{fig:same} compares each 2D baseline with HBM
using the same kind of source and one source per data qubit (Ded). With factories, HBM-Ded
has a lower STV than Planar-2D on 10 of 11 circuits, by 1.31$\times$ on average (geometric
mean; 0.92--1.67$\times$). With cultivation, HBM-Cult-Ded is lower than 2D-Cult on all 11
circuits, by 1.19--1.21$\times$ (1.20$\times$ on average). Every run uses the smallest odd
code distance that keeps its program error below 1\%, so both sides of each ratio meet the
same target.

\textbf{Best layout per circuit:} Across all 13 configurations (Fig.~\ref{fig:heatmap}), the
lowest STV belongs to HBM-Ded on 9 circuits and to HBM-Shared-2 on Ising-10 and Ising-26.
At its best sharing ratio, HBM with factories is below Planar-2D on all 11 circuits
(1.02--1.67$\times$).

\textbf{Denser in-plane cultivation:} Packing the cultivation patches into existing routing
tiles (2D-Cult\textonehalf, one patch per two data qubits) does not close the gap: HBM-Cult-Ded
is 1.29--1.98$\times$ below it.

\begin{rqsummary}
\textbf{RQ1 Summary:} At the same error target and with the same T-state source, an upper
layer lowers STV by 1.31$\times$ with factories and by 1.20$\times$ with cultivation
(geometric means). With its best layout, HBM beats the 2D baseline of its source type on all
11 circuits.
\end{rqsummary}

% ---- RQ2 ------------------------------------------------------------------------------------

\subsection{(RQ2) Qubits or time}

Fig.~\ref{fig:split} splits each ratio of Fig.~\ref{fig:same} into a qubit ratio and a round
ratio.

\textbf{Cultivation, all in qubits:} 2D-Cult and HBM-Cult-Ded take the same number of rounds
on every circuit. Both have a patch next to every data qubit with the same recharge time, so
they run the same schedule. The whole gain is in qubits (1.19--1.21$\times$): in the plane,
the patches add a row to the grid for every row of data qubits; on top, they leave the data
plane unchanged.

\textbf{Factories, qubits for small circuits and rounds for large ones:} On the four
circuits with $N \le 16$, the gain is in qubits (1.31$\times$) at almost equal rounds
(1.00--1.03$\times$): the ring of ports costs area, but its ports keep up with the T gates.
On every larger circuit but Ising-26, Planar-2D takes 1.22--1.84$\times$ more rounds. Its
factories never run out of T states; the extra rounds are T gates whose route to the ring is
blocked by other routes (5,376 such stalls on Heis-64). HBM-Ded serves every T gate from
directly above and runs every circuit at its ideal depth. From $N = 49$ on, it pays for this
in qubits (0.70--0.84$\times$): it builds one factory per data slot, 81 on Ising-66 against
the 40 ports of Planar-2D.

\textbf{Where HBM loses:} On Ising-26, the ring keeps up (1.01$\times$ the rounds) and
HBM-Ded builds 36 factories for 26 data qubits, so it needs 1.10$\times$ the qubits and its
STV is 0.92$\times$ that of Planar-2D. Shared-2 brings it to parity (1.02$\times$).

\begin{rqsummary}
\textbf{RQ2 Summary:} With cultivation, the gain is entirely in qubits (1.20$\times$, same
rounds). With factories, it is in qubits on small circuits (1.31$\times$) and in rounds on
larger ones (1.22--1.84$\times$), where T-gate routes to the perimeter congest; HBM-Ded runs
every circuit at its ideal depth.
\end{rqsummary}

% ---- RQ3 ------------------------------------------------------------------------------------

\subsection{(RQ3) How many sources}

Fig.~\ref{fig:sharing} varies HBM's sharing ratio against the same 2D baselines.

\textbf{Factories:} Shared-2 uses 0.40--1.10$\times$ as many factories as Planar-2D and
roughly matches it (1.05$\times$ on average; lower STV on 6 of 11 circuits). Sparser layouts
lose: 0.61$\times$ at Shared-4 (lower on one circuit), 0.28$\times$ at Shared-8 and
0.12$\times$ at Shared-16.

\textbf{Cultivation:} Every shared layout loses to 2D-Cult, already at Shared-2
(0.59$\times$). A patch delivers at most one T state every 6 ($R = 60$) or 29 ($R = 360$)
steps, so with fewer patches T gates wait.

\textbf{Caveats:} Factories are modelled as always ready, the most favourable case for
sharing. % TODO: factory production time (HBM_FACTORY_ROUNDS=43), pending.
For $N < 16$, Shared-16 is a single source.

\begin{rqsummary}
\textbf{RQ3 Summary:} HBM needs a dense upper layer: one source per data qubit gives the
lowest STV on 9 of 11 circuits, and one per two data qubits on the other two. Sparser
layouts lose to the 2D baselines: with factories from Shared-4 on, and with cultivation
already at Shared-2.
\end{rqsummary}

% ---- figures ----------------------------------------------------------------------------------

\begin{figure}[t]
  \centering
  \includegraphics[width=\columnwidth]{figs_v2/ratios_same.pdf}
  \caption{STV of a 2D baseline divided by that of HBM with the same kind of T source and a
    layout fixed for all circuits, per circuit and as the geometric mean (GMean): Planar-2D
    against HBM-Ded, both with factories, and 2D-Cult (one cultivation patch per data qubit)
    against HBM-Cult-Ded, both with cultivation patches. Ded: one T source per data qubit.
    Above the dashed line (1), HBM has the lower STV. All runs at the smallest odd code
    distance that meets the program-error target, chosen per run ($d = 11$--17). *: the
    circuit's T count requires the slower cultivation grade ($R = 360$ rounds).}
  \label{fig:same}
\end{figure}

\begin{figure}[t]
  \centering
  \includegraphics[width=\columnwidth]{figs_v2/split.pdf}
  \caption{The STV ratios of Fig.~\ref{fig:same} split into their two factors: physical
    qubits ($x$) and rounds ($y$) of the 2D baseline divided by those of HBM, one marker per
    circuit. Circles: Planar-2D over HBM-Ded (factories); squares: 2D-Cult over HBM-Cult-Ded
    (cultivation). Beyond 1 on an axis, HBM uses fewer qubits or rounds; above the dashed
    curve ($xy = 1$), HBM has the lower STV. Markers that coincide share one label.}
  \label{fig:split}
\end{figure}

\begin{figure*}[t]
  \centering
  \includegraphics[width=\textwidth]{figs_v2/heatmap.pdf}
  \caption{STV of every configuration (columns) on every circuit (rows), divided by the
    lowest STV on that circuit; the boxed cell (1.00) is the circuit's best configuration.
    Columns: the 2D baselines (2D-Cult with one or half a cultivation patch per data qubit,
    Planar-2D), HBM with factories on top and HBM-Cult with cultivation patches on top, each
    with one source per data qubit (Ded) or one source shared by 2 to 16 data qubits
    (S2--S16). Row labels give the circuit's T count.}
  \label{fig:heatmap}
\end{figure*}

\begin{figure}[t]
  \centering
  \includegraphics[width=\columnwidth]{figs_v2/sharing.pdf}
  \caption{STV of a 2D baseline divided by that of HBM with the same kind of T source,
    against HBM's sharing ratio (Ded: one T source per data qubit; S2--S16: one source per 2
    to 16 data qubits): Planar-2D over HBM, both with factories (solid), and 2D-Cult (one
    patch per data qubit) over HBM-Cult, both with cultivation patches (dashed). Lines:
    geometric mean over the circuits; bands: range. The two lines use different baselines,
    so their heights are not comparable with each other. Above the dashed line (1), HBM has
    the lower STV; log scale.}
  \label{fig:sharing}
\end{figure}
